\chapter{Diseases of the Ears}
\label{ch:ears}
The ear is the sensory organ responsible for hearing and balance. It is divided into the outer ear (pinna and external auditory canal), middle ear (tympanic membrane, ossicles, and Eustachian tube), and inner ear (cochlea and vestibular apparatus). Diseases of the ear range from common infections to complex disorders of hearing and equilibrium.

%-------------------------------------------------------------
\section{Inner Ear Diseases}
%-------------------------------------------------------------

%-------------------------------------------------------------

%-------------------------------------------------------------

Inner ear diseases involve pathology of the cochlear and vestibular apparatus located within the bony labyrinth of the temporal bone. Disorders of these intricate structures disrupt auditory transduction and spatial orientation mechanisms, leading to sensorineural hearing loss, vertigo, and dysequilibrium. Etiologies include endolymphatic hydrops, autoimmune inflammatory processes, viral infections, vascular compromise, and ototoxic injury.

%-------------------------------------------------------------

%-------------------------------------------------------------

%-------------------------------------------------------------

%-------------------------------------------------------------

\label{sec:Inner Ear Diseases}
%-------------------------------------------------------------%-------------------------------------------------------------

\subsection{Acoustic Neuroma (Vestibular Schwannoma)}
\label{sec:Acoustic Neuroma (Vestibular Schwannoma)}
Vestibular schwannoma is a benign tumor arising from the Schwann cells of the vestibular portion of cranial nerve VIII and is the most common tumor of the cerebellopontine angle. The condition results from neoplastic proliferation of Schwann cells. Bilateral tumors are pathognomonic of neurofibromatosis type 2. Clinical features include unilateral sensorineural hearing loss, tinnitus, and imbalance; larger tumors may compress adjacent structures causing facial numbness (CN V), facial weakness (CN VII), and brainstem compression. Diagnosis is based on MRI with gadolinium, which is the gold standard. Management options include observation with serial imaging (for small, slow-growing tumors), stereotactic radiosurgery, and microsurgical resection. Prognosis is favorable for small tumors, with hearing preservation as a key goal.

\subsection{Benign Paroxysmal Positional Vertigo}
\label{sec:Benign Paroxysmal Positional Vertigo}
Benign paroxysmal positional vertigo is the most common cause of peripheral vertigo, resulting from displaced otoliths from the utricle into the semicircular canals, most commonly the posterior canal. The condition results from movement of calcium carbonate crystals within the semicircular canals. Clinical features include brief episodes (less than 1 minute) of severe vertigo triggered by head movement, with characteristic nystagmus. Diagnosis is based on the Dix-Hallpike test, which reproduces vertigo and reveals torsional-upbeating nystagmus. Management involves canalith repositioning maneuvers (Epley maneuver for posterior canal BPPV), which are highly effective. Prognosis is excellent, although recurrence occurs in 20--50\% of cases.

\subsection{Labyrinthitis}
\label{sec:Labyrinthitis}
Labyrinthitis is an inflammation of the membranous labyrinth of the inner ear that causes both vestibular and auditory dysfunction. The condition results from viral infection, often following an upper respiratory infection; bacterial labyrinthitis is more severe and can result from otogenic or meningogenic spread. Clinical features include vertigo, nausea, vomiting, hearing loss, and tinnitus. Diagnosis is clinical, supported by audiometry and vestibular testing. Management of viral labyrinthitis involves supportive care (vestibular suppressants, antiemetics, followed by vestibular rehabilitation); bacterial labyrinthitis requires aggressive antibiotic therapy and may require labyrinthectomy. Prognosis is favorable for viral labyrinthitis; bacterial labyrinthitis carries a risk of profound hearing loss.

\subsection{Ménière's Disease}
\label{sec:Ménière's Disease}
Ménière's disease is a chronic inner ear disorder characterized by episodic vertigo, fluctuating sensorineural hearing loss (typically low-frequency), tinnitus, and aural fullness. The condition results from endolymphatic hydrops, although the exact pathogenesis remains debated. Onset is typically unilateral, with progression to bilateral disease in 10--30\% of cases over decades. Clinical features include vertigo episodes lasting minutes to hours that may be incapacitating, associated with nausea and vomiting, along with fluctuating hearing loss, tinnitus, and aural fullness. Diagnosis is clinical, supported by audiometry, vestibular testing, and MRI to exclude other causes. Management includes dietary modification (low sodium, caffeine restriction), betahistine, diuretics, intratympanic gentamicin or corticosteroid for refractory cases, and labyrinthectomy or vestibular nerve section as a last resort. Prognosis is variable, with most patients achieving symptom control with treatment.

\subsection{Sudden Sensorineural Hearing Loss}
\label{sec:Sudden Sensorineural Hearing Loss}
Sudden sensorineural hearing loss is an otologic emergency defined as a loss of at least 30 dB over three contiguous audiometric frequencies within 72 hours. The condition results from idiopathic causes, although viral cochleitis, vascular occlusion, autoimmune inner ear disease, or perilymphatic fistula may be implicated. Clinical features include abrupt unilateral hearing loss, often accompanied by tinnitus, vertigo, and aural fullness. Diagnosis is based on audiometry. Management involves high-dose oral or intratympanic corticosteroids, with prompt treatment within 2 weeks being crucial; hyperbaric oxygen therapy may provide additional benefit. Prognosis is variable, with some cases recovering spontaneously and others sustaining permanent hearing loss.

\subsection{Vestibular Neuritis}
\label{sec:Vestibular Neuritis}
Vestibular neuritis is an acute, isolated episode of severe vertigo caused by inflammation of the vestibular nerve without hearing loss. The condition is thought to be viral in etiology, possibly reactivation of herpes simplex virus. Clinical features include sudden-onset, continuous vertigo lasting days to weeks with severe imbalance and nystagmus beating away from the affected side. Diagnosis is based on clinical examination with the head impulse test showing reduced vestibulo-ocular reflex on the affected side. Management involves short-term vestibular suppressants (benzodiazepines, antihistamines) followed by early vestibular rehabilitation. Prognosis is favorable, with most patients recovering fully over weeks.

%-------------------------------------------------------------
\section{Middle Ear Diseases}
%-------------------------------------------------------------

%-------------------------------------------------------------

%-------------------------------------------------------------

%-------------------------------------------------------------

%-------------------------------------------------------------

%-------------------------------------------------------------

%-------------------------------------------------------------

\label{sec:Middle Ear Diseases}
%-------------------------------------------------------------%-------------------------------------------------------------

\subsection{Acute Otitis Media}
\label{sec:Acute Otitis Media}
Acute otitis media is a bacterial or viral infection of the middle ear that affects 80\% of children by age 3. The condition results from Eustachian tube dysfunction secondary to upper respiratory infection, allowing bacterial or viral entry; the most common pathogens are \textit{Streptococcus pneumoniae}, \textit{Haemophilus influenzae}, and \textit{Moraxella catarrhalis}. Clinical features include ear pain, fever, irritability, and hearing loss; otoscopic examination reveals a bulging, erythematous tympanic membrane with loss of landmarks and possible purulent effusion. Diagnosis is based on clinical and otoscopic examination. Management includes analgesics and observation in mild cases, with amoxicillin as first-line antibiotic therapy for moderate-to-severe disease or treatment failure; recurrent acute otitis media may warrant myringotomy with tympanostomy tube placement. Prognosis is excellent with appropriate treatment.

\subsection{Chronic Otitis Media}
\label{sec:Chronic Otitis Media}
Chronic otitis media is a chronic or recurrent perforation of the tympanic membrane with or without active middle ear infection. The condition results from persistent mucosal inflammation or squamous disease (cholesteatoma), which is an abnormal accumulation of keratinizing squamous epithelium that can erode the ossicles, facial nerve, and bony labyrinth. Clinical features include persistent or recurrent otorrhea through a perforated tympanic membrane, hearing loss, and occasional vertigo; cholesteatoma may cause conductive hearing loss, facial nerve palsy, and intracranial complications. Diagnosis is based on otoscopic examination and temporal bone CT. Management is primarily surgical (tympanoplasty, mastoidectomy). Prognosis depends on the extent of disease and success of surgical intervention.

\subsection{Otosclerosis}
\label{sec:Otosclerosis}
Otosclerosis is a primary bone remodeling disorder of the otic capsule that causes progressive conductive hearing loss, typically bilateral, beginning in the third to fourth decade of life, with a female predominance and genetic component (autosomal dominant with variable penetrance). The condition results from abnormal spongy bone formation (active otosclerosis) that replaces normal endochondral bone, most commonly fixing the footplate of the stapes in the oval window. Clinical features include progressive conductive hearing loss. Diagnosis is suggested by a positive Schwartze sign (red hue of the promontory on otoscopy) and confirmed by audiometry showing conductive or mixed hearing loss. Management involves stapedectomy or stapedotomy with prosthetic stapes replacement, which is highly effective; sodium fluoride may slow disease progression. Prognosis is excellent with surgical intervention.

\subsection{Tympanic Membrane Perforation}
\label{sec:Tympanic Membrane Perforation}
Tympanic membrane perforation is a tear in the eardrum that may result from acute otitis media, trauma (penetrating injury, barotrauma, blast injury), or prior tympanostomy tubes. The condition results from mechanical or infectious disruption of the tympanic membrane. Clinical features include hearing loss, brief otalgia, tinnitus, and otorrhea. Diagnosis is based on otoscopic examination. Management involves observation for small perforations, which often heal spontaneously; larger or persistent perforations may require surgical repair (myringoplasty or tympanoplasty). Prognosis is good, with healing enhanced by keeping the ear dry and avoiding nose blowing.

%-------------------------------------------------------------
\section{Outer Ear Diseases}
%-------------------------------------------------------------

%-------------------------------------------------------------

%-------------------------------------------------------------

Outer ear diseases encompass pathology affecting the pinna (auricle) and the external auditory canal extending to the tympanic membrane. These conditions frequently result from mechanical trauma, environmental exposure, infectious pathogens, or dermatological hypersensitivity reactions. Inflammatory and obstructive disorders of the external canal impair sound wave transmission and cause localized pain, otorrhea, and conductive hearing loss.

%-------------------------------------------------------------

%-------------------------------------------------------------

%-------------------------------------------------------------

%-------------------------------------------------------------

\label{sec:Outer Ear Diseases}
%-------------------------------------------------------------%-------------------------------------------------------------

\subsection{Cerumen Impaction}
\label{sec:Cerumen Impaction}
Cerumen impaction is an accumulation of earwax that causes symptoms or obscures the tympanic membrane, affecting approximately 6\% of the general population and up to 30\% of the elderly. The condition results from overproduction of cerumen or from use of cotton swabs and hearing aids that push cerumen deeper into the canal. Clinical features include hearing loss, aural fullness, tinnitus, and otalgia. Diagnosis is based on otoscopic examination. Management involves manual extraction with a curette, irrigation with warm water, or cerumenolytic agents (dilute hydrogen peroxide, mineral oil, glycerin). Prognosis is excellent with proper removal; avoidance of cotton swabs is recommended for prevention.

\subsection{Chondrodermatitis Nodularis Helicis}
\label{sec:Chondrodermatitis Nodularis Helicis}
Chondrodermatitis nodularis helicis is a painful, inflammatory condition affecting the cartilage of the pinna, most commonly the helix or antihelix. The condition results from necrotizing chondritis with overlying epidermal changes, often triggered by chronic pressure during sleep. Clinical features include a small, exquisitely tender nodule with overlying crusting or ulceration, frequently on the side most affected by sleeping pressure. Diagnosis is based on clinical examination and histopathology. Management includes pressure relief (using a donut-shaped pillow), topical corticosteroids, and intralesional corticosteroid injection; refractory cases may require surgical excision. Prognosis is generally good with appropriate treatment.

\subsection{Otitis Externa}
\label{sec:Otitis Externa}
Otitis externa is an infection of the external auditory canal that commonly affects swimmers and individuals with predisposing factors such as water exposure, cotton-swab use, hearing aids, and humid climates. The condition results from inflammation and infection of the canal skin, most commonly caused by \textit{Pseudomonas aeruginosa} and \textit{Staphylococcus aureus}. Clinical features include ear pain (otalgia), pruritus, otorrhea, and a sensation of ear fullness; examination reveals erythema, edema, and tenderness of the canal wall with purulent debris. In severe cases, infection may spread to the surrounding skin (malignant otitis externa), particularly in diabetics and immunocompromised patients. Diagnosis is based on clinical examination. Management involves thorough canal debridement and topical antimicrobial drops (fluoroquinolones or aminoglycosides with corticosteroids); oral antibiotics are reserved for extensive or invasive disease. Prognosis is excellent with appropriate treatment.

%-------------------------------------------------------------
\section{Tinnitus}
%-------------------------------------------------------------

%-------------------------------------------------------------

%-------------------------------------------------------------

%-------------------------------------------------------------

%-------------------------------------------------------------

%-------------------------------------------------------------

%-------------------------------------------------------------

\label{sec:section:Tinnitus}
Tinnitus is the perception of sound in the absence of an external auditory stimulus, affecting millions of individuals worldwide. It is classified as subjective (audible only to the patient, resulting from aberrant neural activity) or objective (audible to an observer, caused by vascular or muscular turbulence). Management focuses on identifying underlying neuro-otological etiologies, sound therapy, cognitive behavioral interventions, and neurostimulation.

%-------------------------------------------------------------%-------------------------------------------------------------

\subsection{Tinnitus}
\label{sec:Tinnitus}
Tinnitus is the perception of sound in the absence of an external auditory stimulus, affecting approximately 10--15\% of adults, with higher prevalence in the elderly. The condition results from maladaptive neural plasticity in the central auditory pathways following peripheral auditory deafferentation; causes include noise-induced hearing loss (the most common cause), age-related hearing loss, ototoxic medications, M\'{e}ni\`{e}re's disease, vestibular schwannoma, otosclerosis, impacted cerumen, and temporomandibular joint disorders. Clinical features include ringing, buzzing, hissing, roaring, or clicking in one or both ears; tinnitus may be associated with depression, anxiety, and insomnia. Diagnosis is based on thorough history, otoscopic examination, and audiometry to assess for hearing loss; MRI with gadolinium is indicated for asymmetric tinnitus or focal neurological findings. Management includes patient education and reassurance, treatment of reversible causes, hearing aids for those with hearing loss, cognitive-behavioral therapy, sound therapy, and tinnitus retraining therapy. Prognosis is variable, with many patients achieving habituation through appropriate management.

